% ======================================================================
% 第 V 节：弱耦合近似
% ======================================================================
\section{弱耦合近似}\label{sec:weak}

弱耦合近似将作简要讨论，留下许多悬而未决的问题。这里只讨论纯规范场。再次考虑
表达式
\begin{equation}
Z_{\theta}(P) = \prod_{n,\mu\nu}\int\ddth{\mu\nu}
  \exp\!\Bigl[i\sum_{P}(\pm)\thl{\mu\nu}
    +\frac{1}{2g^{2}}\sum_{n,\mu\nu}
      \bigl(1-\tfrac{1}{2}\bigl(e^{if_{\mu\nu}(n)}+e^{-if_{\mu\nu}(n)}\bigr)\bigr)\Bigr].
\label{eq:5.1}
\end{equation}

假设积分变量是 $f_{\mu\nu}$ 而非 $\thl{\mu\nu}$。对小 $g$，为使
$\exp[(1/2g^{2})\cos f_{\mu\nu}]$ 接近其最大值 1，积分中只有 $f_{\mu\nu}$ 的小值
才重要。于是可以展开：
\begin{equation}
\frac{1}{2g^{2}}\bigl(1-\cos f_{\mu\nu}\bigr)
  \approx \frac{1}{4g^{2}}f_{\mu\nu}^{2}+\cdots .
\label{eq:5.2}
\end{equation}
用这一近似可以把对 $f_{\mu\nu}$ 的积分限从 $\pm\pi$ 延拓到 $\pm\infty$，误差可忽略；
这样就得到一组高斯积分需要求值。

实际上积分变量是 $\thl{\mu\nu}$，而不是 $f_{\mu\nu}$。不过，可以作从
$\thl{\mu\nu}$ 到 $f_{\mu\nu}$ 的变量变换。这一变换不可能消去全部 $\thl{\mu\nu}$，
而且并非所有变量 $f_{\mu\nu}$ 都独立。尽管如此，该变换足以使 $I_{\theta}(P)$ 对
小的 $g^{2}$ 可计算。

为使变量变换精确，考虑一个有限大小（$|n_{0}|,|n_{1}|,|n_{2}|,|n_{3}| \le N$）
带周期性边界条件的系统。于是可以把变量从 $\thl{\mu\nu}$ 变换到 $f_{\mu\nu}$ 的
一个子集，再加上一些规范变换变量 $\varphi_{n}$，以及四个额外的变量
$\xi_{\mu\nu}$，如下所示：

\emph{(i)} 对满足 $n_{0} < N$ 的 $n$，$\thl{\mu\nu}$（$\mu=1,2,3$）替换为
$f_{\mu\nu}$。对 $\thl{0\mu}$ 写
\begin{equation}
\thl{0\mu} \to f_{0\mu},
\label{eq:5.3}
\end{equation}
并把 $\thl{0\mu}$ 替换为 $\varphi_{\mu}$。这就是对 $\mu \ne 0$ 从 $\thl{\mu\nu}$
到 $f_{\mu\nu}$、对 $\thl{0\mu}$ 到 $\varphi_{\mu}$ 的变换要点。要完成变换，必须
讨论曲面 $n_{0}=N$。

\emph{(ii)} 对 $n_{0}=N$、$n_{1}<N$、$n_{2}$ 和 $n_{3}$ 任意，$\thl{\mu\nu}
(\mu,\nu\in\{1,2,3\})$ 替换为 $f_{\mu\nu}$；$\thl{0\mu}$ 替换为 $\varphi_{\mu}$，
并满足
\begin{equation}
\thl{0\mu}(n) = \varphi_{\mu}(n) + i\,\text{（$f_{\mu\nu}$ 的线性组合）}.
\label{eq:5.4}
\end{equation}

\emph{(iii)} 对 $n_{0}=n_{1}=N$、$n_{2}<N$、$n_{3}$ 任意，
$\thl{\mu\nu}(\mu,\nu\in\{2,3\})$ 替换为 $f_{\mu\nu}$；$\thl{0\mu}$ 替换为
$\varphi_{\mu}$，并满足 $\thl{0\mu}=f_{0\mu}-\varphi_{\mu}$。

\emph{(iv)} 对 $n_{0}=n_{1}=n_{2}=N$、$n_{3}<N$，$\thl{\mu\nu}(\mu=3)$ 替换为
$f_{3\nu}$；$\thl{0\mu}$ 替换为 $\varphi_{\mu}$，其中
$\thl{0\mu}=\varphi_{\mu}-\varphi_{\mu}$。

\emph{(v)} 对 $n_{0}=n_{1}=n_{2}=n_{3}=N$，写 $\thl{\mu\nu}=\xi_{\mu\nu}$，
$\xi_{\mu\nu}$ 是新变量。对该 $n$ 值还令 $\varphi_{n}=0$。

不是积分变量的 $f_{\mu\nu}$（例如对满足 $n_{0}<N$ 的 $n$，$\mu\ne0$ 且
$\nu\ne0$ 的 $f_{\mu\nu}$）可以用式~\eqref{eq:4.6} 表示为独立的 $f_{\mu\nu}$
变量；$\varphi_{n}$ 和 $\xi_{\mu\nu}$ 都不出现在这些表达式中。

当用新的独立变量表示任意的 $\thl{\mu\nu}$ 时，发现（$\xi_{\mu\nu}$ 只在
$n_{0}=N$ 时出现）：
\begin{equation}
\begin{split}
\thl{\mu\nu}(n) = f_{\mu\nu}(n) &+ \text{（规范变换项）}\\
 &+ \text{（$f_{\mu\nu}$ 的线性组合）},
\end{split}
\label{eq:5.5}
\end{equation}
即 $\varphi$ 变量定义了一个规范变换，$f$ 表示对某些 $\theta$ 的平移。容易验证
$Z_{\theta}(P)$ 的被积函数只涉及 $f$：它既与 $\varphi$ 无关，也与 $\xi_{\mu\nu}$
无关（后者不出现，因为任何闭合路径 $P$ 在子格点 $n_{0}=N$ 上具有同样多的
$-\thl{\mu\nu}$ 项和 $+\thl{\mu\nu}$ 项）。因此 $\varphi_{n}$ 与 $\xi_{\mu\nu}$
的积分可以平凡地算出。由于约束式~\eqref{eq:4.6}，这些积分并非平凡。

要完成的目标是：对小 $g$，把格点理论化约成类似常规自由规范场理论的东西。这意味
着恢复 $\thl{\mu\nu}$ 作为积分变量，但积分限为无穷，并加入规范固定项。例如，
假设从下式出发
\begin{equation}
\begin{split}
Z_{\theta}'(P) = \prod_{n,\mu\nu}\int\ddth{\mu\nu}\,
  \exp\!\Bigl[i\sum_{P}(\pm)\thl{\mu\nu}
    &+\frac{1}{2g^{2}}\sum_{n,\mu\nu}
      \bigl(1-\tfrac{1}{2}\bigl(e^{if_{\mu\nu}(n)}+e^{-if_{\mu\nu}(n)}\bigr)\bigr)\\
    &-\frac{\alpha}{2}\sum_{n}\bigl(\nabla_{\mu}A_{\mu}(n)\bigr)^{2}\Bigr],
\end{split}
\label{eq:5.6}
\end{equation}
其中 $\alpha$ 项是 $(\nabla_{\mu}A_{\mu})^{2}$ 规范固定项的格点版本。用显式高斯
积分方法计算这一积分，比计算 $I_{\theta}(P)$ 的 $f_{\mu\nu}$ 积分容易得多。
此外，$I_{\theta}(P)$ 也可以用与 $I_{\theta}(P)$ 相同的变量变换化约为对
$f_{\mu\nu}$ 一个子集的积分。由于 $\alpha$ 项把 $\varphi$ 与 $f$ 耦合，结果在
这种情况下不同；而且 $\xi$ 积分没有收敛因子。要使 $I_{\theta}'(P)$ 有明确定义
且等于 $I_{\theta}(P)$，必须 (a) 为 $\xi$ 积分放入收敛因子，即
$-\tfrac{1}{2}\alpha(\nabla_{\mu}A_{\mu})^{2}$ 项，以及 (b) 在 $\xi$ 上加一个
二次型，以补偿规范固定项的 $\xi$ 积分结果。作者没有完成这一计算；但由于最终结果
仍然是 $Z_{\theta}(P) = I_{\theta}'(P)$ 是对 $\theta$ 的高斯积分，结果大概会与
第二节报告的常规自由场计算相似。
